\section*{Chapter \ref{ch:b2:quadratic} --- Quadratic Forms}

\begin{solution}{exo:b2:quadratic:1}
$q_1 = (x + 2y)^2 - 3y^2$: rank $2$, signature $(1, 1)$ (a
hyperbola-type form).

$q_2$: complete the square in $x$: $q_2 = (x + y)^2 + y^2 + 2yz +
3z^2 = (x+y)^2 + (y + z)^2 + 2z^2$: rank $3$, signature $(3, 0)$
--- positive definite.
\end{solution}

\begin{solution}{exo:b2:quadratic:2}
Eigenvalues $2$ (on $\operatorname{Vect}(1,1,1)$) and $-1$ (on the
plane $x + y + z = 0$). Orthonormalize: $e_1 =
\frac{1}{\sqrt3}(1,1,1)$; in the plane, Gram--Schmidt on $(1,-1,0),
(1,0,-1)$ gives $e_2 = \frac{1}{\sqrt2}(1,-1,0)$, $e_3 =
\frac{1}{\sqrt6}(1,1,-2)$. Then $P = (e_1\ e_2\ e_3)$ is orthogonal
with $P^{\mathsf T}AP = \operatorname{diag}(2,-1,-1)$.

The form $q = 2xy + 2yz + 2zx$ has matrix $A$: in the rotated
coordinates $q = 2X^2 - Y^2 - Z^2$ --- principal axes; signature
$(1,2)$, agreeing with \cref{ex:b2:quadratic:gaussexample}
(the same form!).
\end{solution}

\begin{solution}{exo:b2:quadratic:3}
$u^{**} = u$: $\langle u^{**}x, y\rangle = \langle x, u^*y\rangle =
\langle ux, y\rangle$ for all $y$. $(uv)^* = v^*u^*$: $\langle
uvx, y\rangle = \langle vx, u^*y\rangle = \langle x,
v^*u^*y\rangle$. Kernel: $y \in \ker u^* \iff \langle x, u^*y
\rangle = 0\ \forall x \iff \langle u(x), y\rangle = 0\ \forall x
\iff y \perp \operatorname{im} u$. Ranks: $\dim\ker u^* = n -
\operatorname{rk} u$ (orthogonal complement), so
$\operatorname{rk} u^* = \operatorname{rk} u$ by rank--nullity ---
the Euclidean avatar of the transpose-rank theorem.
\end{solution}

\begin{solution}{exo:b2:quadratic:4}
Spectral: $A = PDP^{\mathsf T}$, $D$ diagonal with entries
$\lambda_i$ satisfying $\lambda_i^3 = \lambda_i$: $\lambda_i \in
\{-1, 0, 1\}$. Then $A^2 = PD^2P^{\mathsf T}$ with $D^2$ diagonal
of entries $0/1$: $A^2$ is symmetric and idempotent
($(A^2)^2 = A^4 = A\cdot A^3 = A^2$) --- symmetric idempotent
$=$ orthogonal projection (it is the projection onto
$\ker(A^2 - I) = \ker(A-I)\oplus\ker(A+I)$ along $\ker A$, and
these are orthogonal by the spectral theorem).

Generally, $P(A) = P\!\left(\text{diag}\right)$: $P(A)$ has the
same eigenvectors, eigenvalues $P(\lambda_i)$ --- the ``spectral
mapping'' at the diagonalizable level.
\end{solution}

\begin{solution}{exo:b2:quadratic:5}
Closed: $O(n) = g^{-1}(\{I\})$ for the continuous $g(P) =
P^{\mathsf T}P$ (polynomial entries). Bounded: each column of $P
\in O(n)$ is a unit vector, so all entries lie in $\intcc{-1}{1}$.
Closed and bounded in $\mathcal{M}_n(\R) \simeq \R^{n^2}$: compact
(\cref{thm:b2:metric:compactprops}~(2)).

Not connected: $\det$ takes the two values $\pm1$ on $O(n)$, and a
continuous surjection onto $\{-1, 1\}$ splits the space
(\cref{ex:b2:metric:glnr}'s argument).
\end{solution}

\begin{solution}{exo:b2:quadratic:6}
\emph{Existence:} $A = PDP^{\mathsf T}$ with $D =
\operatorname{diag}(\lambda_i)$, $\lambda_i \geq 0$; set $B =
P\sqrt D P^{\mathsf T}$ with $\sqrt D =
\operatorname{diag}(\sqrt{\lambda_i})$: symmetric, positive
semidefinite, $B^2 = A$.

\emph{Uniqueness:} let $B$ be symmetric psd with $B^2 = A$. $B$
commutes with $A$; hence $B$ preserves each eigenspace
$E_\lambda(A)$ (for $Ax = \lambda x$: $A(Bx) = BAx = \lambda Bx$).
On $E_\lambda(A)$, the restriction of $B$ is symmetric psd with
square $\lambda\,\mathrm{id}$; its eigenvalues $\mu$ satisfy $\mu^2
= \lambda$, $\mu \geq 0$: $\mu = \sqrt\lambda$ --- so the
restriction, being diagonalizable with the single eigenvalue
$\sqrt\lambda$, \emph{is} $\sqrt\lambda\,\mathrm{id}$. Since $E =
\bigoplus E_\lambda(A)$, $B$ is determined: $B = \sqrt A$.
\end{solution}

\begin{solution}{exo:b2:quadratic:7}
In an orthonormal eigenbasis, $\norm{Ax}_2^2 = \sum \lambda_i^2
x_i^2 \leq (\max_i \lambda_i^2)\norm x_2^2$, with equality at the
corresponding eigenvector: $\vertiii A_2 = \max\abs{\lambda_i}$.
For the given matrix: eigenvalues $3, -1$
(\cref{ex:b2:quadratic:spectralexample}'s twin):
$\vertiii A_2 = 3$.
\end{solution}

\begin{solution}{exo:b2:quadratic:8}
($\Rightarrow$) The top-left $k \times k$ block $A_k$ is the matrix
of the restriction of the (definite) form to the span of the first
$k$ basis vectors: positive definite, so its eigenvalues are
positive and $\Delta_k = \det A_k > 0$.

($\Leftarrow$) Induction on $n$; $n = 1$ clear. Assume all
$\Delta_k > 0$. By induction, $A_{n-1}$ is positive definite: the
form $q$ restricted to $F = \operatorname{Vect}(e_1, \dots,
e_{n-1})$ is definite. Diagonalize $q|_F$ (Gauss): coordinates
$y_1, \dots, y_{n-1}$ with $q|_F = \sum y_i^2$. In the full space,
completing the square in the last variable,
\[
q = \sum_{i=1}^{n-1} \bigl(y_i + c_i x_n\bigr)^2 + c\,x_n^2
\]
for suitable constants (collect the cross terms into the squares).
The reduction exhibits signature $(n-1 + \epsilon, \cdot)$ with
$\epsilon$ the sign contribution of $c$; and the determinant keeps
the sign of the product of the diagonal coefficients under
congruence ($\det(P^{\mathsf T}AP) = (\det P)^2\det A$): $\Delta_n
> 0$ forces $c > 0$. Hence $q$ is a sum of $n$ squares of
independent forms: positive definite.
\end{solution}

\begin{solution}{exo:b2:quadratic:9}
Let $(e_1, \dots, e_n)$ be an orthonormal eigenbasis for $\lambda_1
\geq \dots \geq \lambda_n$.

\emph{$\lambda_2 \leq$ the min-max:} for any hyperplane $H$, the
$2$-dimensional $V = \operatorname{Vect}(e_1, e_2)$ satisfies
$\dim(H \cap V) \geq 1$ (Grassmann): pick a unit $x \in H \cap V$,
$x = ae_1 + be_2$, $a^2 + b^2 = 1$:
\[
\langle u(x), x\rangle = \lambda_1 a^2 + \lambda_2 b^2 \geq
\lambda_2 :
\]
every hyperplane's max is $\geq \lambda_2$.

\emph{$\geq$:} for $H = e_1^{\perp}$, every unit $x = \sum_{i\geq2}
x_ie_i \in H$ has $\langle u(x), x\rangle = \sum_{i \geq 2}
\lambda_i x_i^2 \leq \lambda_2$, attained at $e_2$: this
hyperplane's max is exactly $\lambda_2$. The min over $H$ is
therefore $\lambda_2$.
\end{solution}
